% TOP_PROBLEM: {"problemId": "problem.gauss-circle-problem", "canonicalTitle": "Gauss Circle Problem", "releaseRank": 260, "primaryDomain": "number_theory_arithmetic_geometry", "primaryDomainLabel": "Number theory & arithmetic geometry", "displayStatus": "open", "releaseStatus": "open"}
% !TeX program = lualatex
% ATTEMPT_STATUS: unresolved
\documentclass[11pt]{article}
\usepackage[margin=25mm]{geometry}
\usepackage{fontspec}
\setmainfont{Latin Modern Roman}
\usepackage{amsmath,amssymb,mathtools}
\usepackage{unicode-math}
\setmathfont{Latin Modern Math}
\usepackage{xurl,hyperref}
\hypersetup{colorlinks=true,urlcolor=blue}
\setcounter{secnumdepth}{0}
\setlength{\parindent}{0pt}
\setlength{\parskip}{5pt}
\setlength{\emergencystretch}{3em}
\begin{document}
\section*{\texorpdfstring{Gauss Circle Problem}{Gauss Circle Problem}}\label{260}
\subsection{Definitions and mathematical statement}
For $r\ge0$, define
\[
 N(r)=\#\{(m,n)\in{\mathbb{Z}}^2:m^2+n^2\le r^2\},\qquad E(r)=N(r)-\pi r^2.
\]
The main selected estimate asks whether, for every ${\varepsilon}>0$, there are constants $C_{\varepsilon},r_{\varepsilon}$ with
$|E(r)|\le C_{\varepsilon} r^{1/2+{\varepsilon}}$ for all $r\ge r_{\varepsilon}$.
The broader request is the sharp order of this lattice-counting error. Radius, not squared radius, is the variable; changing variables changes the numerical exponent. The conjectured estimate does not assert the endpoint bound $O(r^{1/2})$ or a fixed sign of the discrepancy.
\par\smallskip\textit{Formulation and concept references:} \href{https://mathworld.wolfram.com/GausssCircleProblem.html}{[S1]}.

\subsection{Short English statement}
How closely does the number of lattice points in a large disk track its area, and is the error bounded by every power just above the square root of its radius?
\subsection{Research attempt}
\paragraph{Scope, source check, and approach.}
This attempt was prepared by GPT-6 Astra Ultra on 27 September 2026.
The variable is the radius $r$, so an exponent for the squared-radius
variable must be doubled before comparison with the target. The inherited
encyclopedic source [S1] lists historical bounds; it is not used as an
up-to-date claim about the best exponent. The primary paper [E1] contains
later improvements. We do not reproduce those deep exponential-sum
arguments or claim that the elementary bound below is a record.

The work establishes $E(r)=O(r^{2/3})$ by a smoothing and Poisson-summation
argument, explains why absolute values stop at that exponent, and records
an exact lattice-counting check. This is weaker than
$O_\varepsilon(r^{1/2+\varepsilon})$ when $\varepsilon<1/6$.

\paragraph{An exact formula and the elementary boundary estimate.}
Let $m=\lfloor r\rfloor$. Separating the origin, the four axes, and
the open quadrants gives
\[
 N(r)=1+4m+4\sum_{a=1}^{m}\left\lfloor\sqrt{r^2-a^2}\right\rfloor.
                                                               \tag{1}
\]
The convention includes boundary points. For integer $r$, every square
root floor can be computed by integer arithmetic, so this identity does
not require a floating-point comparison near a lattice circle.

Let $Q_z=z+[-1/2,1/2]^2$ for $z\in\mathbb Z^2$, and put
$c=\sqrt2/2$. The union of the squares whose centers lie in $B_r$
has area $N(r)$ and lies in $B_{r+c}$. For $r\ge c$, it contains
$B_{r-c}$: the center of a square containing such a point is at distance
at most $c$ from it. Hence
\[
 \pi(r-c)^2\le N(r)\le\pi(r+c)^2,
 \qquad |E(r)|\le2\pi cr+\pi c^2.                       \tag{2}
\]
This proves an $O(r)$ bound, including the jumps at exact lattice radii.
The target requires cancellation beyond a bound on the number of
boundary squares.

\paragraph{Smoothing with an explicit pointwise sandwich.}
Choose a nonnegative $C^\infty$ function $\phi$ supported in the unit
disk, with integral one, and set
$\phi_\delta(x)=\delta^{-2}\phi(x/\delta)$.
For $0<\delta\le1$ and $r\ge2$, triangle inequalities give
\[
 {\boldsymbol1}_{B_{r-\delta}}*\phi_\delta
       \le {\boldsymbol1}_{B_r}
       \le {\boldsymbol1}_{B_{r+\delta}}*\phi_\delta.      \tag{3}
\]
Indeed the left convolution vanishes outside $B_r$ and never exceeds
one; the right convolution is identically one on $B_r$. Boundaries of
the integration domains have measure zero, so (3) also holds at lattice
points on the circle. Summing over $\mathbb Z^2$ sandwiches $N(r)$
between two smooth lattice sums.

Use the Fourier convention
$\widehat f(\xi)=\int_{\mathbb R^2}f(x)e^{-2\pi i x\cdot\xi}\,dx$.
Poisson summation is applicable to the compactly supported smooth
convolutions in (3), and gives, for $s=r\pm\delta$,
\[
 \sum_{z\in\mathbb Z^2}({\boldsymbol1}_{B_s}*\phi_\delta)(z)
 =\pi s^2+\sum_{k\in\mathbb Z^2\setminus\{0\}}
       \widehat{{\boldsymbol1}_{B_s}}(k)\widehat\phi(\delta k).
                                                               \tag{4}
\]
All sums on the right are absolutely convergent after smoothing.
There is no use of an unregularized Fourier series on its jump set.

\paragraph{The disk Fourier bound.}
For $\xi\ne0$, the divergence theorem, applied in the $\xi/|\xi|$
direction, yields a boundary integral and the estimate
\[
 \left|\widehat{{\boldsymbol1}_{B_s}}(\xi)\right|
           \le C s^{1/2}|\xi|^{-3/2}
              \quad(s|\xi|\ge1).                        \tag{5}
\]
Here is the oscillatory estimate needed to justify its power. Rotate
$\xi$ to the first coordinate. The boundary integral is a constant
multiple of
\[
 \frac{s}{|\xi|}\int_0^{2\pi}
          \cos\theta\,e^{-2\pi i s|\xi|\cos\theta}\,d\theta.
\]
Put $\lambda=2\pi s|\xi|$. Neighborhoods of length
$h=\lambda^{-1/2}$ of $0$ and $\pi$ contribute $O(h)$.
On the complementary intervals, integrate by parts using
$\partial_\theta e^{-i\lambda\cos\theta}
 =i\lambda\sin\theta\,e^{-i\lambda\cos\theta}$.
The endpoint terms are $O(1/(\lambda h))$; the integral of the derivative
of $\cos\theta/\sin\theta$ is $O(1/h)$ and has the same prefactor
$1/\lambda$. Thus the angular integral is $O(\lambda^{-1/2})$,
which proves (5). Fixed bounded $\lambda\ge1$ are absorbed in the
constant. This derivation avoids assuming any cancellation between
different dual lattice points.

Repeated integration by parts for the smooth mollifier gives
$|\widehat\phi(\eta)|\le C_M(1+|\eta|)^{-M}$.
Taking $M=3$ in (4), and grouping lattice points in dyadic annuli,
\[
 \sum_{k\ne0}|k|^{-3/2}(1+\delta|k|)^{-3}
       \le C\delta^{-1/2}.                              \tag{6}
\]
For completeness, an annulus $L\le|k|<2L$ contains $O(L^2)$ points
and contributes at most $CL^{1/2}(1+\delta L)^{-3}$.
The geometric sum for $L\le\delta^{-1}$ is $O(\delta^{-1/2})$;
the tail has a summable factor $L^{-5/2}$ after extracting
$\delta^{-3}$. This proves (6) with constants independent of $r$.

\paragraph{Optimizing the two errors.}
The zero-frequency discrepancy between $\pi(r\pm\delta)^2$ and
$\pi r^2$ is $O(r\delta+\delta^2)$. Equations (3)--(6) therefore give
\[
       |E(r)|\le C\bigl(r\delta+r^{1/2}\delta^{-1/2}+1\bigr).
                                                               \tag{7}
\]
Choose $\delta=r^{-1/3}$. Both nonconstant terms are $r^{2/3}$, proving
\[
                              E(r)=O(r^{2/3}).           \tag{8}
\]
The bounded interval $0\le r\le2$ is harmless by (1).
This is a full proof of a classical weaker estimate, not merely a
heuristic based on smoothing length.

\paragraph{The exact obstruction in this route.}
To get a conjectural error $r^{1/2+\varepsilon}$ from (7), the geometric
term would require roughly $\delta\le r^{-1/2+\varepsilon}$.
At that scale the absolute dual sum is as large as
$r^{3/4-\varepsilon/2}$ under our estimate. Thus optimizing absolute
values cannot reach the requested power. One needs cancellation in
\[
 \sum_{k\ne0}\widehat{{\boldsymbol1}_{B_s}}(k)
                                     \widehat\phi(\delta k),
\]
uniformly in the radius and at a substantially finer smoothing scale.
The oscillation within a single Fourier coefficient, already used in
(5), must not be confused with cancellation among those coefficients.

A second tempting argument treats the fractional parts in (1) as
independent random variables. They are deterministic, depend on the
same $r$, and are constrained by representations as sums of two squares.
A square-root fluctuation calculation for independent random variables
therefore supplies no uniform bound for this sum. Likewise a bound
holding for almost every radius, or an integral second moment, would
not by itself establish the required estimate for every sufficiently
large real radius.

\paragraph{Boundary and computational checks.}
At a radius $r=\sqrt m$, the jump size is the number of representations
$m=a^2+b^2$. Formula (3) remains valid there, so the proof has not silently
removed the most delicate radii. Substituting $x=r^2$ in (8) gives an
$O(x^{1/3})$ bound, while the conjectural radius exponent becomes
$x^{1/4+\varepsilon/2}$; confusing these would double the claimed gain.

The local verification script compares (1) with direct enumeration for
integer radii through a stated finite range and checks (2) numerically
only as a sanity check of constants. The count itself uses integer
square roots. Such computations test implementations and signs, not
the asymptotic conjecture or its endpoint.

\paragraph{Final status.}
\textbf{UNRESOLVED.} The checked partial result is the classical
$O(r^{2/3})$ estimate, with a pointwise smoothing argument valid also
at jump radii. The first missing lemma is a stronger uniform estimate
for the signed nonzero-frequency sum at scales near $r^{-1/2}$.
Specific next targets are grouping equal dual lengths before taking
absolute values, estimating short annuli with their phases intact,
and tracking the dependence on the smoothing scale in such estimates.
No endpoint $O(r^{1/2})$ bound, sharp-order claim, or new exponent is
established here.

\subsection{Sources}
\begin{itemize}
\item[S1]Formal statement, Bounds on a solution and conjecture.
\url{https://mathworld.wolfram.com/GausssCircleProblem.html}.
\textit{Formulation source.}
\item[E1]Jean Bourgain and Nigel Watt, \emph{Mean square of zeta function, circle problem and divisor problem revisited}, 2017.
\url{https://arxiv.org/abs/1709.04340}.
\textit{Primary context for later circle-problem estimates; not a proof input to the elementary bound above.}
\end{itemize}
\end{document}
